\documentclass[12pt]{article}

% --------------------
% PAQUETES BÁSICOS
% --------------------
\usepackage[spanish,es-noshorthands]{babel}
\usepackage[utf8]{inputenc}      % Si usas pdflatex
\usepackage[T1]{fontenc}
\usepackage{lmodern}
%\usepackage[authoryear]{natbib}
\usepackage{comment}
\usepackage{amsmath, amssymb, amsthm}
\usepackage{geometry}
\usepackage{booktabs}   
\usepackage{setspace}
\usepackage{graphicx}
\usepackage{tikz}
\usepackage{tikz-cd}
\usepackage[
  colorlinks=true,
  linkcolor=blue,
  citecolor=blue,
  urlcolor=red
]{hyperref}

\newtheorem{thm}{Teorema}[section]
\newtheorem{dfn}[thm]{Definición}
\newtheorem{lem}[thm]{Lema}
\newtheorem{ej}[thm]{Ejemplo}
\newtheorem{cor}[thm]{Corolario}
\newtheorem{obs}[thm]{Observación}
\newtheorem{con}[thm]{Conjetura}
\newtheorem{prop}[thm]{Proposición}

\DeclareMathOperator{\Frm}{Frm}
\DeclareMathOperator{\Fil}{Fil}
\DeclareMathOperator{\Ord}{Ord}
\DeclareMathOperator{\Top}{Top}
\DeclareMathOperator{\Sp}{sp}
\DeclareMathOperator{\pt}{pt}
\DeclareMathOperator{\id}{id}
\DeclareMathOperator{\tp}{tp}
\DeclareMathOperator{\fd}{fd}
% --------------------
% CONFIGURACIÓN
% --------------------
\geometry{margin=2.5cm}
\onehalfspacing

% --------------------
% DATOS DEL DOCUMENTO
% --------------------
\title{Reporte mensual de actividades}
\date{1-agosto-2026 a 31-agosto-2026}
\author{Juan Carlos Monter Cortés}

% --------------------
% INICIO DEL DOCUMENTO
% --------------------
\begin{document}

\maketitle
% \tableofcontents

\section{Actividades realizadas}

Si consideramos $A\in \Frm$ y $j\in NA$, entonces $A_j$ es un cociente cerrado si $j=u_a$ para algún $a\in A$. Adicionalmente, $A_j$ es un cociente compacto si $\nabla(j)\in A^\wedge$. 
Además, para $F=\nabla(j)$, se satisface que para todo $j\in [v_F, w_F]$, $A_j$ es compacto. Dentro de los objetivos de nuestra investigación se encuentra el estudiar aquellos núcleos que producen 
cociente compactos y cerrados, es decir, aquellos $j\in [v_F, w_F]$ tales que $j=u_a$ para algún $a\in A$ y $F=\nabla(j)\in A^\wedge$. Para realizar dicho estudio, damos la siguiente definición.

\begin{dfn}\label{KC}
Decimos que $A\in \Frm$ es $KC$ si todo cociente compacto es cerrado.
\end{dfn}

Anteriormente, verificamos que para cualquier terna HM $(F, Q, \nabla)$ se satisface que si $A$ es eficiente y $\mathcal{U}\circ v_F=v_\nabla\circ \mathcal{U}$, entonces $A$ es fuertemente apilado.
Haremos uso de este hecho para ver la relación de eficiente y $KC$. 

\begin{comment}

Antes de ello, necesitamos algunos resultados preliminares.\\

Primero, definimos el siguiente subconjunto de $A$
\[
uA=\{a\in A\mid A_{u_a} \text{ es un cociente compacto}\}.
\]

Notemos que a todo $a\in uA$ le asignamas un subconjunto $\nabla(u_a)\in A^\wedge$, es decir, podemos definir la función
\[
\alpha\colon uA\to A^\wedge.
\]

Veamos que propiedades satisface dicha función.

\begin{lem}\label{lem:injetiva}
    Para $uA\subseteq A$ como antes, las siguientes afirmaciones se cumplen:
    \begin{enumerate}
        \item Si $A$ es eficiente, entonces $\alpha$ es suprayectiva.
        \item Si $u_a$ es ajustado para todo $a\in uA$, entonces $\alpha$ es inyectiva.
        \item Si $A$ es eficiente y $a\in uA$, entonces $\alpha$ es biyectiva si y solo si $u_a$ es ajustado.
    \end{enumerate}
\end{lem}

\begin{proof}
$\mbox{ }$
\begin{enumerate}
    \item Consideremos $F=\nabla(j)\in A^\wedge$ y $v_F$ el núcleos ajustado asociado a $F$. Si $A$ es eficiente, entonces $v_F=u_d$,
    donde $d=v_F(0)\in A$. Luego, 
        \[
        \alpha(d)=\nabla(u_d)=\nabla(v_F)=F.
        \]
    Por lo tanto $\alpha$ es suprayectiva.
    \item  Consideremos $a, b\in uA$ tales que $\alpha(a)=\alpha(b)$, es decir, $\nabla(u_a)=\nabla(u_b)$. Por hipótesis, $u_a, u_b$ son núcleos ajustados, es decir, 
    \[
    u_c\leq u_d\quad\mbox{y}\quad u_c\geq u_d.
    \]
    Evaluando ambos núcleos en $0$ obtenemos $c=d$ y por lo tanto $\alpha$ es inyectiva.
    \item Sean $A$ eficiente, $a\in uA$ y supongamos que $\alpha$ es un biyectiva. Consideremos $F=\nabla(u_a)$ y $v_F$ el núcleo ajustado asociado $F$. Por la eficiencia, $v_F=u_d$, es decir,
    \[
    \alpha(a)=\nabla(u_a)=\nabla(v_F)=\nabla(u_d)=\alpha(d)
    \]
    donde $d=v_F(0)$.\\

    Por hipótesis, $\alpha$ es biyectiva, entonces, por la inyectividad, $a=d$. Por lo tanto $v_F=u_a$ y $u_a$ es ajustado.\\

    Reciprocamente, como $A$ es eficiente, entonces por 1), $\alpha$ es suprayectiva. Por hipótesis, $u_a$ es ajustado y, por 2), $\alpha$ es inyectiva. Por lo tanto $\alpha$ es biyectiva. 
\end{enumerate}
\end{proof}

Dado que todo marco $KC$ es eficiente, podemos dar una versión del lema anterior involucrando los marcos $KC$.

\begin{cor}\label{cor:KC}
        Sean $A$ un marco $KC$ y $F\in A^\wedge$. $\alpha$ es biyectiva si y solo si para todo $j\in [v_F, w_F]$, $j$ es ajustado.
\end{cor}

\begin{proof}
Supongamos que $A$ es $KC$ y sea $F\in A^\wedge$. Por definición, para todo $j\in [v_F, w_F]$, $j=u_a$ para algún $a\in A$. Dado que $A$ es $KC$, entonces $A$ es eficiente y por el lema anterior, $\alpha$ es biyectiva si y solo si para todo $j\in [v_F, w_F]$, $j$ es ajustado.
\end{proof}
\end{comment}

\begin{prop}\label{prop:KC}
Sea $A$ un marco. Si $A$ es eficiente y 
\begin{equation}\label{eq:KC}
\mathcal{U}\circ v_F=v_\nabla\circ \mathcal{U},
\end{equation}
entonces $A$ es $KC$.
\end{prop}

\begin{proof}
\begin{comment}
    Supongamos que $A$ es $KC$. Trivialmente, $KC$ implica eficiente y así, resta verificar que $KC$ implica que $\mathcal{U}\circ v_F=v_\nabla\circ \mathcal{U}$ o equivalentemente, por propiedades del adjunto derecho,
    \[
    v_F=\mathcal{U}_*\circ v_\nabla\circ \mathcal{U}.
    \]
    
    Si $A$ es $KC$, entonces para cualquier $F\in A^\wedge$ se cumple que para todo $j\in [v_F, w_F]$, $j=u_a$ para algún $a\in A$. En particular,
    \[
    v_F=u_d,\quad\text{y}\quad \mathcal{U}_*\circ v_\nabla\circ \mathcal{U}=u_a
    \]
    donde $d, a\in A$.\\

    Por construcción, siempre se cumple que $u_d\leq u_a$. Para ver la otra desigualdad verificamos que la función $\alpha$ es inyectiva. Notemos que $d\leq a$  
\end{comment}

Supongamos que $A$ es eficiente y $\mathcal{U}\circ v_F=v_\nabla\circ \mathcal{U}$ para toda terna HM $(F, Q, \nabla)$. Primero, por la eficiencia, 
$A$ es $T_1$. Además, si $A$ es eficiente y $\mathcal{U}\circ v_F=v_\nabla\circ \mathcal{U}$, entonces $A$ es fuertemente apilado. De aquí que, al ser $A$ $T_1$ y fuertemente apilado, 
se cumple que $v_F=w_F$ y, por la eficiencia, $v_F=u_d$, donde $d=v_F(0)$. Por lo tanto, para todo $j\in [v_F, w_F]$, $j=u_d$, es decir, $A$ es $KC$.
\end{proof}

\begin{cor}\label{FapiladoKC}
    Si $A$ es eficiente y fuertemente apilado, entonces $A$ es $KC$.
\end{cor}

Dado que $KC$ implica eficiente, debemos preguntarnos si $KC$ implica la igualdad \eqref{eq:KC}. De manera particular, podriamos investigar si $KC$ implica fuerte apilamiento. Si algunas de las condiciones resulta ser cierta, entonces tendriamos el recíproco de la Proposición \ref{prop:KC} o del Corolario \ref{FapiladoKC}.\\

Algunos de los reportes previos exploran la relación de las propiedades de apilamiento junto con la propiedad $T_1$ en marcos. Dado que en el caso espacial, 
\[
T_1+\text{ apilado}\Longrightarrow \text{sobrio},
\]
al llevar cada propiedad al caso de marcos, la implicación correspondiente es
\[
T_1+\text{ apilado}\Longrightarrow \text{espacial}.
\]

Hemos notado que las proposiciones analogas a las presentadas en \cite{sexton2006ordinal,simmons2004vietoris} donde se involucran $T_1$ y apilado, 
cuando las llevamos al caso de marcos, no implican espacialidad. 
Por lo tanto, nos preguntamos lo siguiente:
\[
\text{¿Existen marcos } T_1 \text{ y apilados que no son espaciales?}
\]

\begin{ej}\label{No espacial}
Sea $B$ un algebra booleana completa sin atomos, entonces $\pt B=\emptyset$. Considerando $\mathbf{2}=\{0<1\}$, definimos
\[
A=B\times \mathbf{2}
\]
el cual es un álgebra booleana completa. Estudiemos algunas de las propiedades de $A$.\\

Primero, veamos que $A$ es un marco $T_1$. Para esto, consideremos $p=(1,0)$ y sean $(a,b), (c,d)\in A$ tales que
\[
(a,b)\wedge (c, d)\leq p.
\]

Debemos verificar que $(a,b)\leq p$ o $(c,d)\leq p$. Dado que el orden en $A$ es coordenada a coordenada, tenemos que
\[
(a\wedge c, b\wedge d)\leq (1,0)\Rightarrow a\wedge c\leq 1 \text{ y } b\wedge d\leq 0.
\]
Como $a\wedge c\leq 1$ siempre es cierta, la condición $b\wedge d\leq 0$ implica que $b=0$ o $d=0$. Por lo tanto, 
\[
(a,b)\leq p \quad \text{o} \quad (c,d)\leq p,
\]
es decir, $p\in \pt A$.\\

Supongamos ahora que existe $q=(x, y)\in \pt A$, entonces $q\neq (1,1)$. Asi, si $y=0$,
\[
(1,0)\wedge (x, 1)\leq (x,0)=q\Rightarrow (1, 0)\leq q\quad\text{o}\quad (x, 1)\leq q,
\]

Ya que $(x, 1)\leq q$ no puede ocurrir, entonces $(1,0)\leq q$, lo cual ocurre si $x=1$. Por lo tanto, $q=(1,0)=p$.\\

Si $y=1$, entonces $x<1$. Luego, consideremos $a, b\in B$ tales que $a\wedge b\leq x$. Entonces
\[
(a,1)\wedge (b,1)\leq (x,1)=q\Rightarrow (a,1)\leq q\quad\text{o}\quad (b,1)\leq q.
\]
Así, $a\leq x$ o $b\leq x$, lo cual implica que $x\in \pt B$. Lo anterior es una contradicción.\\

Por lo tanto, $\pt A=\{p\}$.\\

Ahora, supongamos que existe $(a, b)\in A$ tal que $p\leq (a, b)$. Entonces
\[
(1,0)\leq (a,b)\Rightarrow a=1\; \text{y}\ 0\leq b.
\]
De aquí que, $(a, b)=(1, 0)=p$ o $(a, b)=(1,1)$, es decir, $p$ es máximo y $A$ es un marco $T_1$.\\

El marco $A$ no es espacial, pues si consideramos $a, b\in B$ tales que $a\neq b$, entonces $(a, 0)\neq (b, 0)$ y
\[
\mathcal{U}(a,0)=\{p\in \pt A\mid (a,0)\nleq p\}=\emptyset=\{p\in \pt A\mid (b,0)\nleq p\}=\mathcal{U}(b,0). 
\]
Así, $\mathcal{U}$ no es inyectiva y $A$ no es espacial.\\

Para verificar que $A$ es apilado, hacemos uso de los siguientes hechos:
\begin{enumerate}
    \item En un álgebra booleana completa, los filtros abiertos son principales.
    \item Si $B$ es un álgebra booleana completa sin átomos, entonces $0$ es el único elemento compacto de $B$.
    \item Sea $A$ un marco y $F=\uparrow a$, con $a\in A$. $F$ es un filtro abierto si y solo si $a$ es compacto.
    \item Sea $A$ un marco y $(F, Q)$ un par HM. Si $v_F(0)=w_F(0)$, entonces $A$ es apilado.
\end{enumerate}

Considerando $K(-)$ como el conjunto de elementos compactos de un marco y, por los hechos anteriores, tenemos que
\[
K(A)=K(B\times \mathbf{2})=K(B)\times K(\mathbf{2})=\{(0,0),(0,1)\}.
\]

Luego,
\[
A^\wedge=\{\uparrow (0,0), \uparrow (0,1)\}.
\]

Notemos que $\uparrow (0,0)=A$ y 
\[
\uparrow (0,1)=\{(a,b)\in A\mid (0,1)\leq (a,b)\}=B\times \{1\}.
\]

Además, por la correspondencia biyectiva entre puntos y filtros completamente primos, tenemos que para $p=(1,0)\in \pt A$,
\[
F_p=\{(a,b)\in A\mid (a,b)\nleq p\}=\{(a,1)\in A\mid a\in B\}=\uparrow (0,1).
\]

Por lo tanto, $A^\wedge =\{A, F_p\}$.\\

Dado que $A$ es $T_1$, para el filtro $F_p$, se puede verificar que 
\[
v_{F_p}=u_p=w_p=w_{F_p}.
\]

Así, $v_{F_p}(0)=w_{F_p}(0)$.\\

Cuando $F=A$, entonces 
\[
v_F=\tp=w_F
\]
y $v_F(0)=w_F(0)$.\\

Por lo tanto, para todo $F\in A^\wedge$, se cumple que $v_F(0)=w_F(0)$ y $A$ es apilado.
\end{ej}

De esta manera, el Ejemplo \ref{No espacial} verifica lo siguiente:

\begin{thm}
Existen marcos $T_1$ y apilados que no son espaciales.
\end{thm}

Adicionalmente, en el Ejemplo \ref{No espacial} se satisface que para todo $F\in A^\wedge$, $v_F=w_F$. Lo anterior implica que $A$ es fuertemente apilado, es decir,
\[
\text{existen marcos }T_1 \text{ y fuertemente apilados que no son espaciales.}
\]

En \cite{PasekaSmarda1992} proporcionan una familia de marcos que satisfacen la propiedad $\mathbf{(H)}$ y que no son espaciales. Por lo tanto, ya que 
\[
\mathbf{(H)}\Longrightarrow T_1,
\]
basta con ver que $\mathbf{(H)}$ implica apilado para exhibir otra familia de marcos no espaciales.

\begin{prop}
Si $A\in \Frm$ satisface $\mathbf{(H)}$, entonces $A$ es apilado.
\end{prop}

\begin{proof}
Consideremos $A$ un marco y $(F, Q)$ un par HM. Notemos que si $A$ satisface $\mathbf{(H)}$, entonces $S=\pt A$ es $T_2$ y por lo tanto 
\[
\mathcal{U}_*\circ[Q']\circ \mathcal{U}=w_F
\]
donde $w_F(x)=\bigwedge\{p\in Q\mid x\leq q\}$.

Debemos verificar que, para $s=(\mathcal{U}_*\circ [Q']\circ \mathcal{U})(0)=\bigwedge Q$,
\[
s\leq v_F(0).
\]

A manera de contradicción, supongamos que $s\nleq v_F(0)$. Dado que $w_F\neq \tp$, $1\neq s$ y, por hipótesis, existe $c\in A$ tal que
\[
1)\;c\nleq s\quad\text{y}\quad 2)\;\neg c\nleq v_F(0)
\]

De 1), notemos que existe $r\in Q$ tal que $c\nleq r$, pues en caso contrario, si 
\[
\forall r\in Q\text{ tal que }c\leq r\Rightarrow c\leq \bigwedge Q=s
\]
lo cual es una contradicción.\\

Por 2), si $\neg c\nleq v_F(0)$, entonces $c\notin F$, pues si $c\in F$, 
\[
\neg c=v_c(0)\leq f_F(0)\leq v_F(0)
\]
lo cual sería una contradicción. Así, existe $q\in Q$ tal que $c\leq q$.\\

Por lo tanto, de 1) y 2) obtenemos que existe $c\in A$ y $q, r\in Q$ tales que
\[
c\leq q \quad\text{y}\quad c\nleq r.
\]

Además, por $T_1$, $q\neq r$. Notemos que 
\[
r\vee q=r\vee c=1,
\]
entonces
\[
\neg r\leq c\leq q\quad \text{y}\quad \neg q\leq \neg c\leq r. 
\]

Por HM, sabemos que para todo $a\in F$ se satisface que 
\[
Q\subseteq \mathcal{U}(a)\iff [\forall p\in Q](a\nleq p),
\]
en particular $a\nleq q, r$, es decir, $a\vee q=a\vee r=1$. Por lo tanto,
\[
\neg q, \neg r\leq a 
\]
para todo $a\in F$. Luego, $(\neg q\vee \neg r)\leq a$ para todo $a\in F$.\\

Consideremos $d=\neg q\vee \neg r$. De aquí que
\[
d\leq a\Rightarrow v_a\leq v_d\Rightarrow v_F\leq v_d.
\]

Evaluando en $0$ obtenemos que $v_F(0)\leq \neg d$.
\end{proof}

\begin{comment}

\begin{thm}
Existen marcos $T_1$ y apilados que no son espaciales.
\end{thm}

\begin{proof}
Si consideramos $A\in \Frm$ tal que $A$ satisface $\mathbf{(H)}$, podemos construir el siguiente marco
\[
\widetilde{A}=\{(a,b)\mid a\in A, b\in \mathcal{B}(A)\},
\]
donde $\mathcal{B}(A)=\{b\in A\mid b=\neg \neg b\}$ y se conoce como la \emph{booleanización} de $A$.\\

Además, si para $x\in A$, con $x$ máximo, se cumple que $\neg x=0$, entonces $\widetilde{A}$ satisface $\mathbf{(H)}$. En consecuencia, $\widetilde{A}$ es un marco $T_1$ y apilado, pero no es espacial.
\end{proof}

Para ver los detalles de la no espacialidad de $\widetilde{A}$, se puede consultar \cite{PasekaSmarda1992}.
\end{comment}

%\bibliographystyle{plainnat}
\bibliographystyle{amsalpha}
\bibliography{research2}

\end{document}